\documentclass[border=5mm]{standalone}

\usepackage[T1]{fontenc}
\usepackage[utf8]{inputenc}
\usepackage{tikz}
\usetikzlibrary{arrows.meta}

\begin{document}

\begin{tikzpicture}[
    scale=0.8,
    every node/.style={font=\small\sffamily},
    % Styl dla linii dendrogramu
    dendro line/.style={thick, draw=black},
    % Styl dla osi i siatki
    axis line/.style={thin, draw=black, -{Stealth[scale=0.8]}},
    grid line/.style={help lines, dotted, draw=gray!40}
]

    % --- 1. OŚ PIONOWA (Odległość / Podobieństwo) ---
    \draw[axis line] (-0.8, 0) -- (-0.8, 10.5) node[above, font=\small\bfseries\sffamily] {distance / similarity};

    % Podziałka i linie siatki na osi Y
    \foreach \y/\label in {0/0, 2/20, 4/40, 6/60, 8/80, 10/100} {
        \draw[grid line] (-0.8, \y) -- (9, \y);
        \draw (-0.8, \y) -- (-0.9, \y) node[left] {\label};
    }

    % --- 2. ETYKIETY OBIEKTÓW (Oś X) ---
    \node (A) at (1, 0) [below=3pt] {object A};
    \node (B) at (3, 0) [below=3pt] {object B};
    \node (C) at (5, 0) [below=3pt] {object C};
    \node (D) at (7, 0) [below=3pt] {object D};
    \node (E) at (9, 0) [below=3pt] {object E};

    % --- 3. STRUKTURA DRZEWA (DENDROGRAM) ---

    % Połączenie A i B na wysokości 20 (Y = 2)
    \draw[dendro line] (1,0) -- (1,2) -- (3,2) -- (3,0);
    \coordinate (AB) at (2,2); % Środek gałęzi AB

    % Połączenie D i E na wysokości 30 (Y = 3)
    \draw[dendro line] (7,0) -- (7,3) -- (9,3) -- (9,0);
    \coordinate (DE) at (8,3); % Środek gałęzi DE

    % Połączenie C z grupą (D,E) na wysokości 60 (Y = 6)
    \draw[dendro line] (5,0) -- (5,6) -- (8,6) -- (8,3);
    \coordinate (CDE) at (6.5,6); % Środek gałęzi CDE

    % Ostateczne połączenie grupy (A,B) z grupą (C,D,E) na wysokości 90 (Y = 9)
    \draw[dendro line] (2,2) -- (2,9) -- (6.5,9) -- (6.5,6);


    % --- 4. LINIA ODCIĘCIA (Kryterium podziału) ---
    % Linia na wysokości Y = 5 (odpowiada odległości 50)
    \draw[black, dashed, thick] (-1.2, 5) -- (9.5, 5)
        node[right, font=\footnotesize\bfseries\sffamily, fill=white, inner sep=2pt] {cut-off line};

\end{tikzpicture}

\end{document}
